% !TEX root =./ECE448_Practice_main.tex
%
% L10 practice set (ASCII over PAM). Every question is a step in the L10 link:
% characters to bits, a packet's length and timing, the access-code search,
% what noise and the matched filter do to the text (watched, not calculated:
% ROADMAP, "Sources and noise", 2026-10-21), and the receiver settings that must match
% the transmitter's. Blank copy: \begin{solution}[h] reserves h of work space.
% SOLUTIONS copy (\iskey=1) prints the red worked solutions.

\fullwidth{\textbf{Learning Objective 4.4: } I can implement a PAM link that transmits and
receives ASCII text end-to-end.
\\[4pt]\normalfont\small Show your work. Write a \textbf{Documentation} line at the
end (\textbf{None} if you did not collaborate).}

\begin{questions}

%%%%%%%%%%%%%%%%%%%%%%%%%%%%%%%%%%%%%%%%%
\question \textbf{Text to bits and back.} The L10 transmitter sends each byte most
significant bit first. In ASCII, `H' is 0x48 and `i' is 0x69.
\begin{parts}

	\begin{minipage}[t]{0.58\linewidth}
		\part Write the 16 bits the transmitter sends for ``Hi'', in order.
		\begin{solution}[0.9in]
		0x48 is 01001000 and 0x69 is 01101001, so the bits are
		\texttt{01001000 01101001}.
		\end{solution}
	\end{minipage}\hfill%
	\begin{minipage}[t]{0.37\linewidth}
		\raggedleft \vspace{12pt}
		\ansbox[1.6in]{01001000 01101001}
	\end{minipage}

	\begin{minipage}[t]{0.58\linewidth}
		\part A receiver's Repack Bits block is set to LSB first by mistake. What byte
		does it produce from the first 8 bits, and does it print as a letter?
		\begin{solution}[0.9in]
		Read in reverse, \texttt{01001000} becomes \texttt{00010010}, which is 0x12. That is a
		control code, not a printable character, so the line arrives as nonsense.
		\end{solution}
	\end{minipage}\hfill%
	\begin{minipage}[t]{0.37\linewidth}
		\raggedleft \vspace{12pt}
		\ansbox[1.1in]{0x12, no}
	\end{minipage}

	\begin{minipage}[t]{0.58\linewidth}
		\part A correctly set receiver outputs \texttt{01000111 01001111}. What two
		characters arrived? (`A' is 0x41, and the capital letters follow in order.)
		\begin{solution}[0.9in]
		\texttt{01000111} is 0x47, six after 0x41, which is `G'. \texttt{01001111} is
		0x4F, fourteen after 0x41, which is `O'. The text is ``GO''.
		\end{solution}
	\end{minipage}\hfill%
	\begin{minipage}[t]{0.37\linewidth}
		\raggedleft \vspace{12pt}
		\ansbox[1.1in]{GO}
	\end{minipage}

\end{parts}

\newpage
%%%%%%%%%%%%%%%%%%%%%%%%%%%%%%%%%%%%%%%%%
\question \textbf{One packet.} A cadet changes the L10 message to
\texttt{MAYDAY 41.0N 104.8W}, 19 characters. The Vector Source adds a newline, and the
packet format is the L10 one: a 64-bit access code, a 32-bit header, then the payload.
The link runs at 4 kbaud with one bit per symbol.
\begin{parts}

	\begin{minipage}[t]{0.58\linewidth}
		\part How many bits are in one packet, and how long does it take to send?
		\begin{solution}[1.1in]
		The payload is $19 + 1 = 20$ bytes.
		\eq{ n &= 64 + 32 + 8 \times 20 = 256\ \text{bits} \\
		T &= \dfrac{256}{4000\ \text{baud}} = 64\ \text{ms} }
		\end{solution}
	\end{minipage}\hfill%
	\begin{minipage}[t]{0.37\linewidth}
		\raggedleft \vspace{12pt}
		bits: \ansbox[1.1in]{256}

		\vspace{4pt}
		time: \ansbox[1.1in]{64 ms}
	\end{minipage}

	\begin{minipage}[t]{0.58\linewidth}
		\part Write the 16-bit length field, which the header sends twice.
		\begin{solution}[0.8in]
		The length is 20 bytes, which is \texttt{0000000000010100} in 16 bits.
		\end{solution}
	\end{minipage}\hfill%
	\begin{minipage}[t]{0.37\linewidth}
		\raggedleft \vspace{12pt}
		\ansbox[1.6in]{0000000000010100}
	\end{minipage}

	\begin{minipage}[t]{0.58\linewidth}
		\part What fraction of the bits on the air are text? If the same packets ran at
		9.6 kbaud, the rate of L9's channel example, how long would one take?
		\begin{solution}[1.0in]
		\eq{ \dfrac{160}{256} = 62.5\%, \qquad \dfrac{256}{9600\ \text{baud}} = 26.7\ \text{ms} }
		\end{solution}
	\end{minipage}\hfill%
	\begin{minipage}[t]{0.37\linewidth}
		\raggedleft \vspace{12pt}
		text: \ansbox[1.1in]{62.5\%}

		\vspace{4pt}
		time: \ansbox[1.1in]{26.7 ms}
	\end{minipage}

\end{parts}

\newpage
%%%%%%%%%%%%%%%%%%%%%%%%%%%%%%%%%%%%%%%%%
\question \textbf{The access-code search.} Correlate Access Code - Tag Stream compares
the last 64 bits with the access code after every new bit and declares a packet when at
most 2 bits differ.
\begin{parts}

	\part While no packet is arriving, the bits are random. About how many of the 64
	match the access code at a typical position, and why?
		\begin{solution}[0.8in]
		About 32 bits match. Each random bit matches the code's bit with probability one half, so on
		average half of the 64 match.
		\end{solution}

	\part A teammate sets the threshold to 0, so the block declares a packet only when
	all 64 bits match. In words, what happens to the link as the noise rises, compared
	with a threshold of 2? Why not set the threshold to 20 instead?
		\begin{solution}[1.3in]
		With a threshold of 0, a single wrong bit anywhere in the access code loses the
		whole packet, so lines stop arriving at a lower noise setting even when the
		payload would have come through intact. A threshold of 2 lets a packet with one
		or two damaged access-code bits through. A threshold of 20 is far too loose:
		random bits match about 32 of 64 by chance, and some stretches come within 20,
		so the receiver would start false packets in the middle of the noise and data.
		\end{solution}

	\part On the PlutoSDR version of the link, no packets arrive, yet the constellation
	after the Costas Loop shows two tight clusters on the real axis. Explain in words what the
	correlator sees and why, and what the receiver should change.
		\begin{solution}[1.2in]
		The signal arrives inverted: the carrier-recovery loop has locked 180$^\circ$ off, so
		every 1 arrives as a 0. Where the access code passes, every bit differs, so the
		correlator counts 0 of 64 matches instead of 64, and it never declares a packet.
		The fix is to flip the polarity, multiplying the samples by $-1$ before the
		slicer, which is what the lab flowgraph's polarity switch does.
		\end{solution}

\end{parts}

\newpage
%%%%%%%%%%%%%%%%%%%%%%%%%%%%%%%%%%%%%%%%%
\question \textbf{Watching noise break the text.} The L10 link runs at a noise
amplitude of 0.3 with every line intact. A cadet raises the noise slider a step at a
time while watching the matched-filter eye, the Symbol Sync output, and the Received and
Lines labels.
\begin{parts}

	\part Describe, in order, what changes on each display as the noise rises from 0.3
	to about 1.4.
		\begin{solution}[1.4in]
		The eye fills in from above and below until it closes. The Symbol Sync output,
		two tight bands at $\pm 1$, spreads into two clouds that reach toward zero. Then
		single characters go wrong in the Received label, the correct count falls behind
		the received count, and finally the received count falls too, because a wrong bit
		in the access code or header loses the whole line. By about 1.4, no line arrives
		intact.
		\end{solution}

	\begin{minipage}[t]{0.58\linewidth}
		\part At 1.1, the line reads \texttt{USAFA ECE 448: ASCII over PCM.} instead of
		\texttt{\ldots PAM.} `A' is 0x41 and `C' is 0x43. How many bits went wrong in
		that character, and which one?
		\begin{solution}[1.0in]
		0x41 is \texttt{01000001} and 0x43 is \texttt{01000011}: one bit, the second
		from the right, flipped from 0 to 1.
		\end{solution}
	\end{minipage}\hfill%
	\begin{minipage}[t]{0.37\linewidth}
		\raggedleft \vspace{12pt}
		\ansbox[1.1in]{1 bit}
	\end{minipage}

	\part A teammate proposes 200-byte packets to raise the share of text on the air
	from 72\% to 94\%. In words, would the longer line start to break at a lower or a
	higher noise setting than the 31-byte line, and why?
		\begin{solution}[1.1in]
		Lower. A line arrives intact only if every one of its bits is right, and a 200-byte
		line has about six times as many bits that can go wrong, so at any noise setting
		it is more likely to catch at least one wrong bit.
		\end{solution}

	\part At a noise amplitude of 0.5, every line arrives. The cadet switches the
	matched filter off, and only a few lines in a hundred arrive intact. Explain why,
	and say which display shows it first.
		\begin{solution}[1.1in]
		The signal occupies only a few kilohertz of the 32 kHz the samples represent. The
		matched filter blocks the noise outside the signal's band; without it, all that
		noise reaches Symbol Sync and the slicer, and many samples cross zero. The eye
		shows it first: its opening disappears the moment the filter is switched off.
		\end{solution}

\end{parts}

\newpage
%%%%%%%%%%%%%%%%%%%%%%%%%%%%%%%%%%%%%%%%%
\question \textbf{Clock error.} Two radios each derive their clocks from a crystal
accurate to 20 ppm, so the receiver's sample clock can differ from the transmitter's
symbol clock by up to 40 ppm. The link runs at 4 kbaud.
\begin{parts}

	\begin{minipage}[t]{0.58\linewidth}
		\part A receiver without timing recovery samples at fixed intervals. How many
		symbols, and how many seconds, until its sampling point slips one whole symbol?
		How many 344-bit packets is that?
		\begin{solution}[1.3in]
		\eq{ N &= \dfrac{1}{40\times10^{-6}} = 25{,}000\ \text{symbols} \\
		t &= \dfrac{25{,}000}{4000} = 6.25\ \text{s} \\
		\dfrac{25{,}000}{344} &= 73\ \text{packets} }
		\end{solution}
	\end{minipage}\hfill%
	\begin{minipage}[t]{0.37\linewidth}
		\raggedleft \vspace{12pt}
		time: \ansbox[1.1in]{6.25 s}

		\vspace{4pt}
		packets: \ansbox[1.1in]{73}
	\end{minipage}

	\part The L10 link loses almost no lines up to 10,000 ppm with Symbol Sync in the
	receiver. Explain why the slip in part (a) does not appear in that link, and why
	a 40 ppm error is not a concern for it.
		\begin{solution}[1.1in]
		Symbol Sync measures the timing error from the samples near each zero crossing and
		moves its sampling instant to follow the drift, so the sampling point never
		slips. 40 ppm is 250 times smaller than the 10,000 ppm the link handles,
		so the loop follows it.
		\end{solution}

	\part In the clock-error screenshot, the eye sink looks smeared while every line
	arrives. Explain why the eye sink and Symbol Sync disagree.
		\begin{solution}[1.0in]
		The eye sink cuts the waveform every 8 samples, the nominal samples per symbol, so
		when each symbol spans slightly fewer samples, each trace starts a little later in
		its symbol than the last and the eye drifts across the display. Symbol Sync
		follows the actual symbol timing, so its samples stay at the eye's center.
		\end{solution}

\end{parts}

\vspace{1.5em}
\noindent\textbf{Documentation:} \rule[-2pt]{0.6\linewidth}{0.4pt}

\end{questions}
